\subsection{Finitely generated abelian groups}

In the case \(\mathbf{A} = \mathbf{Z}\) , the results of section 4 can be expressed:

THEOREM 3. --- Every finitely generated abelian group G is a direct sum of its torsion subgroup F (subgroup of elements of finite order in G) and a free abelian group of finite rank p (isomorphic to \(Z^{p}\) ). The group F is a direct sum of a finite number of cyclic groups of orders \(n_{1}\) , \(n_{2}\) , \ldots, \(n_{q}\) , where the \(n_{i}\) are integers \textgreater1, each of which divides the previous one; moreover, the integers p, q and \(n_{i}\) ( \(1 \leq i \leq q\) ) are uniquely determined by G.

Remark. --- While the orders \(n_{1}, \ldots, n_{q}\) of the cyclic groups of which F is the direct sum are well defined by the divisibility condition of Th. 3, it is not the same for the groups themselves: for example, in the product G of \(\mathbf{Z}/(p)\) with itself (for p prime), the subgroups are precisely the \(F_{p}\) -vector subspaces, and G is the direct sum of two 1-dimensional subspaces in \(p(p+1)\) different ways.

COROLLARY 1. --- In a finite abelian group G, there exists an element whose order is the lcm of all the orders of elements of G; this order n is the first invariant factor of G.

COROLLARY 2. --- Any finite abelian group G whose order is not divisible by the square of any integer \textgreater1 is cyclic.

Let us keep to the notation of Th. 3. Then \(p = 0\) since \(G\) is finite, and \(q \leqslant 1\) , for otherwise the order of \(G\) would be divisible by \(n_q^2\) . Hence \(G\) is cyclic.

COROLLARY 3. --- Let L and M be two free Z-modules of rank n, let \((e_{i})\) be a basis for L and \((f_{i})\) a basis for M \((1 \leqslant i \leqslant n)\) , let u be a homomorphism from L into M, and let U be its matrix with respect to the bases \((e_{i})\) and \((f_{i})\) . Then Coker \(u = \mathrm{M}/u(\mathrm{L})\) is finite if and only if \(\det(U) \neq 0\) , and then \(\operatorname{Card}(\operatorname{Coker}(u)) = |\det(U)|\) .

By changing bases in L and M if necessary, we may assume that U has the form described in VII, p.~22, Cor. 1 to Prop. 5 (where the \(\alpha_{i}\) are in this case integers); the corollary then becomes obvious, since the order of a direct sum of Z-modules \(Z/\alpha_{i}Z\) ( \(1 \leqslant i \leqslant n\) ) is infinite if one of the \(\alpha_{i}\) is zero, and is equal to \(|\alpha_{1}\alpha_{2} \ldots \alpha_{n}|\) otherwise.

